\documentclass[DIV=calc]{scrartcl}
\usepackage[utf8]{inputenc}
\usepackage[T1]{fontenc}
\usepackage[ngerman]{babel}
\usepackage{graphicx}

\usepackage{fixltx2e}
\usepackage{ellipsis}
\usepackage[tracking=true]{microtype}

\usepackage{lmodern}                        % Ersatz fuer Computer Modern-Schriften
\usepackage{hfoldsty}

\usepackage{fourier} 			% Schriftart
\usepackage[scaled=0.81]{helvet} 	% Schriftart

\usepackage{url}
\usepackage{tocloft} 			% Paket für Table of Contents

\usepackage{xcolor}
\definecolor{urlred}{HTML}{660000}

\usepackage{hyperref}
\hypersetup{
  colorlinks=true,	
  linkcolor=black,	% Farbe der internen Links (u.a. Table of Contents)
  urlcolor=black,	% Farbe der url-links
  citecolor=black} % Farbe der Literaturverzeichnis-Links

\usepackage{mdwlist} 	% Änderung der Zeilenabstände bei itemize und enumerate

\parindent 0pt 				% Absatzeinrücken verhindern
\parskip 12pt 				% Absätze durch Lücke trennen

\usepackage{titlesec}	% Abstand nach Überschriften neu definieren
\titlespacing{\subsection}{0ex}{3ex}{-1ex}
\titlespacing{\subsubsection}{0ex}{3ex}{-1ex}		

% \pagestyle{empty}
\setlength{\textheight}{22.5cm}
\usepackage{fancyhdr}
\pagestyle{fancy}
\cfoot{}
\lfoot{Zusammenkunft aller Physik-Fachschaften\\43,0te Konferenz der Informatikfachschaften\\Konferenz der deutschsprachigen Mathematikfachschaften}
\rfoot{zapfev.de\\kif.fsinf.de\\die-koma.org}
\renewcommand{\headrulewidth}{0pt}
\renewcommand{\footrulewidth}{0.1pt}

\usepackage{adjustbox}

\begin{document}
\begin{adjustbox}{center}
\begin{minipage}{1.2\textwidth}
	\begin{minipage}{4cm}
		\vspace{-1.4cm}
		\includegraphics[width=\textwidth]{komaLogo.png}\\
		\centering
		\small Konferenz der \\ Mathematikfachschaften
	\end{minipage}
	\hfill
	\begin{minipage}[c]{0.25\textwidth}
		\vspace{-2cm}
		\includegraphics[width=\textwidth]{zkk_logo.pdf}
	\end{minipage}
	\hfill
	\begin{minipage}{4cm}
		\begin{minipage}{120pt}
			\vspace{-1.8cm}
			\includegraphics[width=80pt]{../../logo.pdf}
			\centering
			\small Zusammenkunft aller Physik-Fachschaften
		\end{minipage}
	\end{minipage}
\end{minipage}
\end{adjustbox}
\vspace{0.3cm}
\begin{center}
\huge{Stellungnahme der Zusammenkunft aller Physik-Fachschaften, der Konferenz
der Informatikfachschaften 43,0 und der Konfrenz der deutschsprachigen Mathematikfachschaften} \\
\normalsize
\end{center}

\vspace{0.4cm}
\section*{Resolution zur qualitativen Umsetzung von  eduroam und anderen
hochschulöffentlichen Netzwerken an allen Hochschulen}

Die zeichnenden Bundesfachschaftentagungen begrüßen das weit verbreitete
Angebot von eduroam an deutschsprachigen Hochschulen und halten die
Qualitätssicherung des eduroam-Netzwerkes für die Arbeit an Hochschulen für
unerlässlich.

Wir halten folgende Punkte für besonders kritisch und möchten daher auf diese
explizit hinweisen:

\begin{enumerate}
\item Wir fordern die Einhaltung der eduroam Policy Service Definition,
  festgelegt von der GÉANT Association, in der Version 2.8 vom Juli 2012, da
  in der Vergangenheit von einigen Hochschulen einige Empfehlungen sowie
  Forderungen hierin nicht beachtet wurden.

  Herausheben wollen wir dabei
  \begin{enumerate}
  \item  die Einhaltung der in Abschnitt 6.3.3, Unterpunkt ``Network'',
    aufgeführten Liste der unbedingt anzubietenden Ports. Leider wurden wir
    auf zahlreiche Verstöße gegen diesen Punkt aufmerksam gemacht.

    Wir unterstützen darüber hinaus die Empfehlung, keine
    bzw. möglichst wenige Portrestriktionen vorzunehmen, sowie keine
    Anwendungs- und Abfangproxies zu verwenden.

  \item die Einhaltung der in Abschnitt 6.3.2 festgelegten Unterstützung
    anonymer Authentifizierung. Wir bitten, diese Unterstützung auch in den
    entsprechenden Anleitungen zu dokumentieren.
  \end{enumerate}
\item Falls Portrestriktionen unumgänglich sind, sollten diese öffentlich
  zugänglich dokumentiert und begründet werden, sowohl für ein- als auch für
  ausgehende Beschränkungen.

  Wir bitten die GÉANT Association, dies in die eduroam Policy Service Definition
  als ``MUST''-Anforderung aufzunehmen.

\item Aufgrund der herausragenden Bedeutung des eduroam-Netzes für die
  wissenschaftliche Gemeinschaft fordern wir eine ausreichende Ausstattung mit
  personellen und finanziellen Mitteln zur Aufrechterhaltung, zur Verbesserung
  und zum Ausbau des Netzwerkes durch die betreibenden Hochschulen.
\end{enumerate}

Wir bitten diese Hinweise analog für andere hochschulöffentliche Netze zu
beherzigen.

\vfill
\begin{flushright}
Verabschiedet durch die KIF 43,0 und die KoMa am 30.05.2015 in Aachen \\
Verabschiedet durch die ZaPF am 31.05.2015 in Aachen
\end{flushright}




\end{document}
